Human-written and AI-written text are linearly separable in the residual stream of language models, and recent detectors exploit this. It is not known whether such a direction also \emph{controls} how AI-like a model's own writing sounds, or whether it is simply formality, verbosity, domain, or the Assistant persona under another name.
We extract a content-matched difference-of-means direction \dai from \qwen and test it causally. We steer generation along \dai and score the outputs with independent detectors, while controlling for coherence and content and comparing against equal-norm random, confound, and prompting baselines.
Steering toward ``human'' at block 20 lowers \desklib \pai from 0.996 to 0.84 and raises the zero-shot \binoculars score from 0.70 to 0.83, while three random directions leave \desklib at 0.99 (paired difference $-0.15$, $p<10^{-10}$). The effect holds among generations judged perfectly fluent (0.85 \versus 0.98--1.00 for formality and random directions). It also survives projecting out six confound directions.
\dai is almost orthogonal to the Assistant axis in whitened geometry (cosine 0.005). It nearly coincides with an instruct-minus-base direction from another model family, and it transfers to chat generators but not to base LMs. It is already present in the base model (cosine 0.99 with the instruct direction) and steers it in both directions.
The lever is weak, however. Fluent steered text stays far from human answers (0.85 \versus 0.35), and an LLM judge finds formality steering and prompting equally effective.
``Sounding like AI'' is thus a real, causal, chat-tuning style feature, but one direction does not capture what makes text read as AI.
